\documentclass[11pt,reqno]{amsart}

\usepackage[T1]{fontenc}
\usepackage{lmodern}
\usepackage{microtype}
\usepackage{amsmath,amssymb,amsthm,mathtools}
\usepackage{xcolor}
\usepackage[
  colorlinks=true,linkcolor=blue!45!black,citecolor=blue!45!black,
  urlcolor=blue!45!black,
  pdftitle={Real arithmetic schedules and operation counts for fastCD},
  pdfauthor={Harrison Lemley}
]{hyperref}

\numberwithin{equation}{section}
\newtheorem{proposition}{Proposition}[section]
\newtheorem{corollary}[proposition]{Corollary}

\newcommand{\R}{\mathbb R}
\newcommand{\C}{\mathbb C}
\newcommand{\Rk}{\operatorname{R}}
\newcommand{\ReC}{\operatorname{Re}_{\mathrm c}}
\newcommand{\ImC}{\operatorname{Im}_{\mathrm c}}
\newcommand{\AltC}{\mathsf C}
\newcommand{\AltR}{\mathsf R}

\title[Real arithmetic schedules for fastCD]{Real arithmetic schedules
and operation counts for fastCD}
\author{Harrison Lemley}
\date{October 2026}
\thanks{The measured source snapshot is preserved with the benchmark data
at fastCD revision
\href{https://git.tanuki-cd.com/algebraity/fastCD/commit/%
4954189686d70c4419ee669abdf9caefe3818dd4}{\texttt{4954189686}},
including the archived
\href{https://git.tanuki-cd.com/algebraity/fastCD/src/commit/%
4954189686d70c4419ee669abdf9caefe3818dd4/%
benchmark-results/paper-v2-octonion243-20261023/source.patch}{source patch}.}

\begin{document}
\maketitle

\section{Notation}
\label{sec:notation}

The notation and results in this technical note are based on the mathematical
framework and results of
\cite{Lemley2026}.

Let $R=\R$ and $S=\C$, and fix nonzero real doubling parameters
$\mu_k$ for $k\ge1$. Set $A_0:=\R$ and define
$A_n:=A_{n-1}\oplus A_{n-1}$ with conjugation and multiplication
\begin{equation}\label{eq:CD-product}
\overline{(a,b)}=(\bar a,-b)\,,\qquad
(a,b)(c,d)=\bigl(ac+\mu_n\bar d\,b,\;da+b\bar c\bigr)\,,
\end{equation}
where $\bar r=r$ for $r\in\R$. Write $N=2^n$ and use the recursive
standard basis $e_0=1,e_1,\ldots,e_{N-1}$. The standard real tower has
$\mu_k=-1$ for every $k$. Put
\[
V_n:=\{x\in A_n:\bar x=-x\}\,,\qquad
A_n^{\C}:=A_n\otimes_{\R}\C\,,\qquad
V_n^{\C}:=V_n\otimes_{\R}\C\,.
\]
Thus $V_n$ consists of the elements whose $e_0$ coordinate is zero.
Define the bilinear pairing in these coordinates by
\begin{equation}\label{eq:weights}
\langle x,y\rangle_n:=\sum_{j=0}^{N-1}w_{n,j}x_jy_j\,,\qquad
w_0=(1)\,,\qquad w_n=(w_{n-1},-\mu_n w_{n-1})\,.
\end{equation}
Conjugation is extended $\C$-linearly and the pairing $\C$-bilinearly.
For $X,Y\in V_n$, write
\[
X\times_nY:=\frac12(XY-YX)\,,
\]
and let $\AltC_n$ denote its $\C$-bilinear extension to $V_n^{\C}$.
For $x=\alpha+X$ and $y=\beta+Y$, with $\alpha,\beta\in\R$ and
$X,Y\in V_n$, full multiplication is recovered by
\begin{equation}\label{eq:scalar--imaginary}
xy=\bigl(\alpha\beta-\langle X,Y\rangle_n\bigr)
   +\alpha Y+\beta X+X\times_nY\,.
\end{equation}

At level $n\ge2$, write $X=(u,\gamma+s)$ and $Y=(v,\delta+t)$, where
$u,v,s,t\in V_{n-1}^{\C}$ and $\gamma,\delta\in\C$, and put
\[
\begin{aligned}
D&:=\AltC_{n-1}(u,v)+\mu_n\AltC_{n-1}(s,t)\,,\\
H&:=\AltC_{n-1}(u,t)+\AltC_{n-1}(s,v)\,,\\
\lambda&:=\langle s,v\rangle_{n-1}-\langle t,u\rangle_{n-1}\,.
\end{aligned}
\]
Expanding \eqref{eq:CD-product} gives
\begin{equation}\label{eq:doubled-alt}
\AltC_n(X,Y)=\bigl(D+\mu_n(\delta s-\gamma t),\;
                    \lambda+\delta u-\gamma v-H\bigr)\,.
\end{equation}
The notation $B_n^{\C}(X,Y)$ below uses the triples
$(u,\gamma,s)$ and $(v,\delta,t)$ from this splitting. For real
coefficients, define the nonrecursive bilinear map by
\begin{equation}\label{eq:nonrecursive-map}
\begin{aligned}
B_n\bigl((u,\gamma,s),(v,\delta,t)\bigr)
:=\bigl(&\mu_n(\delta s-\gamma t),\;
\langle s,v\rangle_{n-1}-\langle t,u\rangle_{n-1}\,,\\
&\delta u-\gamma v\bigr)\,.
\end{aligned}
\end{equation}
Write $B_n^{\C}$ for its $\C$-bilinear extension and put
$(F,\lambda,G)=B_n^{\C}(X,Y)$. For $n\ge3$, choose
$\rho_n\in\C$ with $\rho_n^2=\mu_n$ and compute
\[
Z_\pm:=\AltC_{n-1}(u\pm\rho_n s,v\pm\rho_n t)\,.
\]
The transformed inputs belong to $V_{n-1}^{\C}$ because that space is
closed under complex linear combinations. Expanding by bilinearity gives
\[
\begin{aligned}
Z_\pm
&=\AltC_{n-1}(u,v)
 \pm\rho_n\AltC_{n-1}(u,t)
 \pm\rho_n\AltC_{n-1}(s,v)
 +\rho_n^2\AltC_{n-1}(s,t)\\
&=D\pm\rho_n H\,.
\end{aligned}
\]
Since $\mu_n\ne0$, also $\rho_n\ne0$, and hence
\[
D=\frac{Z_++Z_-}{2}\,,\qquad
H=\frac{Z_+-Z_-}{2\rho_n}\,.
\]
Substitution into \eqref{eq:doubled-alt} gives the recursive coefficient
kernel
\begin{equation}\label{eq:two-call-formula}
\AltC_n(X,Y)=\Bigl(\frac{Z_++Z_-}{2}+F,\;
                   \lambda+G-\frac{Z_+-Z_-}{2\rho_n}\Bigr)\,.
\end{equation}
Its level-one base case is
$\AltC_1\equiv0$, and its level-two base case is the direct formula
$(F,\lambda+G)$, with the scalar $\lambda$ inserted in the unit coordinate
of the second component. The standard-tower schedule uses the octonion formula
below
as its level-three base case. For real inputs,
write $\AltR_n(X,Y)=X\times_nY$ and use the real-input reductions of
Section~\ref{subsec:real-specialization}, followed by
\eqref{eq:scalar--imaginary} for full multiplication.

Write $T_{A_n}$ for the real multiplication tensor. Its real tensor rank
$\Rk_{\R}(T_{A_n})$ is the least $r$ for which
\[
xy=\sum_{j=1}^r\alpha_j(x)\beta_j(y)v_j
\]
with real linear forms $\alpha_j,\beta_j$ and fixed $v_j\in A_n$.
In this bilinear model, fixed linear coordinate changes have no cost.


\section{Real coefficient specialization}
\label{subsec:real-specialization}

Write $\iota$
for the central coefficient satisfying $\iota^2=-1$. The Cayley--Dickson
involution is
extended $\C$-linearly and therefore fixes $\iota$, whereas coefficientwise
complex conjugation sends $\iota$ to $-\iota$. Let $\ReC$ and $\ImC$ denote
the real and imaginary coefficient projections. Complex addition uses $2$
real additions, multiplication by $\iota$ requires only a coordinate exchange
and a sign change, and complex multiplication can be expressed using $3$ real
products:
\begin{equation}\label{eq:three-real-complex}
\begin{gathered}
(a+\iota b)(c+\iota d)=(p-q)+\iota(r-p-q)\,,\\
p=ac\,,\qquad q=bd\,,\qquad r=(a+b)(c+d)\,.
\end{gathered}
\end{equation}

Write $X=a_0+\iota a_1$ and $Y=b_0+\iota b_1$, where
$a_0,a_1,b_0,b_1\in V_n$. Applying this multiplication formula to the
nonrecursive map $B=B_n$ in \eqref{eq:nonrecursive-map} gives
\begin{equation}\label{eq:complex-fusion}
\begin{gathered}
\Delta=B(a_0,b_0)\,,\qquad \Omega=B(a_1,b_1)\,,\\
W=B(a_0+a_1,b_0+b_1)\,,\\
B^{\C}(a_0+\iota a_1,b_0+\iota b_1)
=(\Delta-\Omega)+\iota(W-\Delta-\Omega)\,.
\end{gathered}
\end{equation}
Thus the complex kernel evaluates its nonrecursive map by three real evaluations
of $B$, followed by reconstruction of the complex output coefficients.

\begin{corollary}[One-call real-input root]\label{cor:real-root}
Let $n\ge3$ and $\mu_n<0$. Put
$a_n=\sqrt{-\mu_n}$ and $\rho_n=a_n\iota$. For
$X=(u,\gamma+s)$ and $Y=(v,\delta+t)$ in $V_n$, define
\[
Z:=\AltC_{n-1}(u+\rho_n s,v+\rho_n t)\,,\qquad
\lambda:=\langle s,v\rangle_{n-1}-\langle t,u\rangle_{n-1}\,.
\]
Then
\begin{equation}\label{eq:real-root}
X\times_nY
=\Bigl(\ReC Z+\mu_n(\delta s-\gamma t),\;
\lambda+\delta u-\gamma v-a_n^{-1}\ImC Z\Bigr)\,.
\end{equation}
\end{corollary}

\begin{proof}
For $X,Y\in V_n$, the vectors $u,v,s,t$ have real coefficients. Hence the
quantities
\[
D=\AltC_{n-1}(u,v)+\mu_n\AltC_{n-1}(s,t)
\]
and
\[
H=\AltC_{n-1}(u,t)+\AltC_{n-1}(s,v)
\]
also have real coefficients. Expanding $Z$ by bilinearity and using
$\rho_n=a_n\iota$ gives
\[
Z=D+a_n\iota H\,.
\]
Therefore
\[
\ReC Z=D\,,\qquad\ImC Z=a_nH\,.
\]
Substituting these identities into \eqref{eq:doubled-alt} gives
\eqref{eq:real-root}.
\end{proof}

This one-call reduction applies at the real-input root when $\mu_n<0$,
while general complex recursive calls still use the two-call formula.
For $\mu_n>0$, choosing $\rho_n=\sqrt{\mu_n}\in\R$ gives real transformed
inputs, so two real-input calls at the preceding level suffice. The real-input
schedule uses these reductions for $n\ge3$.

For Propositions~\ref{prop:octonion-leaf}--\ref{prop:exact-counts}, take
$\mu_k=-1$ for every $k$, with $\rho_k=\iota$,
so the tower is standard.
Write $M$ for input-dependent real multiplications, $A$ for real additions
and subtractions, and $E$ for multiplications by fixed scalars other than
$0,\pm1$. Each costs $1$ operation, so the total is $T=M+A+E$.
Known-zero terms are omitted and sums are initialized from their first term.
Preprocessing of both variable operands is included in the total cost, while
coordinate extraction, copying, permutations, and sign changes are
not counted, and scalar and imaginary parts are obtained directly from
the standard-basis coordinates.

\section{An octonion base case}
\label{subsec:octonion-leaf}

When $n=3$, the complex kernel uses the following formula for the
alternating product $\AltC_3$ on the complexified imaginary octonions,
whose elements have $7$ complex coordinates. With the one-call real-input
root, this base case first occurs in full multiplication in $A_4$, of real
dimension $16$.

Write an input to $\AltC_3$ in the standard coordinates as
\[
X=(x_1,x_2,x_3,h,y_1,y_2,y_3)\in V_3^{\C}\,.
\]
For $u,v\in\C^3$, let
\[
u\times v:=(u_2v_3-u_3v_2,\;u_3v_1-u_1v_3,\;u_1v_2-u_2v_1)
\]
be the ordinary three-dimensional cross product, extended bilinearly to
complex coordinates. Define the following linear forms in $X$:
\[
\begin{gathered}
p=-\iota h\,,\qquad z=x_1+p\,,\\
b=x_2+\iota y_2\,,\qquad r=x_2-\iota y_2\,,\qquad
c=x_3+\iota y_3\,,\qquad s=x_3-\iota y_3\,.
\end{gathered}
\]
Let $L_j\colon\C^7\to\C^3$, $1\le j\le4$, be the linear maps
\begin{equation}\label{eq:octonion-linear-maps}
\begin{aligned}
L_1X&=(x_1,y_1,p)\,,
& L_2X&=(b+s,r-c,p)\,,\\
L_3X&=(z+\iota y_1,b,c)\,,
& L_4X&=(z-\iota y_1,r,s)\,.
\end{aligned}
\end{equation}
Their transposes below are ordinary matrix transposes, without complex
conjugation.

\begin{proposition}[Octonion leaf]\label{prop:octonion-leaf}
For $X,Y\in V_3^{\C}$,
\begin{equation}\label{eq:octonion-leaf}
\AltC_3(X,Y)
=-\iota L_1^{\mathsf T}(L_1X\times L_1Y)
 +\frac12\sum_{j=2}^4L_j^{\mathsf T}(L_jX\times L_jY)\,.
\end{equation}
Under the arithmetic convention above,
this formula has an evaluation using $72$ input-dependent real
multiplications, $156$ real additions and subtractions, and $15$ real
halvings, for a total of $243$ operations.
\end{proposition}

\begin{proof}
To verify the identity, express the alternating trilinear form associated
with $\AltC_3$ as a sum of determinants. For complex linear forms
$f,g,\ell$ on $\C^7$, write $[f\,g\,\ell](X,Y,Z)$ for the determinant
whose columns are $(f(X),g(X),\ell(X))$, $(f(Y),g(Y),\ell(Y))$, and
$(f(Z),g(Z),\ell(Z))$. The bracket is multilinear and alternating in
$f,g,\ell$ and in $X,Y,Z$, and its arguments will be suppressed below.
In particular, expansion in the first two forms gives
\[
\begin{aligned}
[f+\iota g\,\ f-\iota g\,\ \ell]
&=-\iota[f\,g\,\ell]+\iota[g\,f\,\ell]\\
&=-2\iota[f\,g\,\ell]\,.
\end{aligned}
\]
Let
\[
\Phi(X,Y,Z):=\langle\AltC_3(X,Y),Z\rangle_3\,.
\]
Since the parameters are standard, the pairing in this expression is
the complex-bilinear dot product. In the standard coordinates of
$V_2^{\C}$, the alternating product is the ordinary cross product and
the pairing is $u\cdot v:=\sum_{j=1}^3u_jv_j$. Write
$X=(x,h,y)$, $Y=(x',h',y')$, and $Z=(x'',h'',y'')$, where
$x=(x_1,x_2,x_3)$ and $y=(y_1,y_2,y_3)$, with primes indicating the
corresponding coordinates of $Y$ and $Z$. Formula~\eqref{eq:doubled-alt}
gives
\[
\begin{aligned}
\AltC_3(X,Y)=\Bigl(&x\times x'-y\times y'+hy'-h'y,\;\\
&y\cdot x'-y'\cdot x,\;
h'x-hx'-x\times y'-y\times x'\Bigr)\,.
\end{aligned}
\]
Pairing with $Z$ therefore yields
\[
\begin{aligned}
\Phi={}&(x\times x'-y\times y'+hy'-h'y)\cdot x''\\
&+(y\cdot x'-y'\cdot x)h''\\
&+(h'x-hx'-x\times y'-y\times x')\cdot y''\,.
\end{aligned}
\]
For each index $j$, the terms involving $h,h',h''$ are
\[
\begin{aligned}
&(hy_j'-h'y_j)x_j''
 +(y_jx_j'-y_j'x_j)h''
 +(h'x_j-hx_j')y_j''\\
&\qquad=[x_j\,h\,y_j](X,Y,Z)\,.
\end{aligned}
\]
The remaining dot and cross products expand into the four determinants
in the first line below. Consequently,
\begin{equation}\label{eq:octonion-three-form}
\begin{aligned}
\Phi={}&[x_1\,x_2\,x_3]-[x_1\,y_2\,y_3]
       +[x_2\,y_1\,y_3]-[x_3\,y_1\,y_2]\\
&+[x_1\,h\,y_1]+[x_2\,h\,y_2]+[x_3\,h\,y_3]\,.
\end{aligned}
\end{equation}
Here the coordinate symbols in brackets denote the coordinate linear
forms, rather than their values at a particular input.

Put $a=x_1+\iota y_1$ and $q=x_1-\iota y_1$. Expanding
$[b+s\,\ r-c\,\ p]$ gives
\[
[b\,r\,p]-[b\,c\,p]+[s\,r\,p]-[s\,c\,p]\,.
\]
The term $-[b\,c\,p]$ cancels $[p\,b\,c]$ from
$[a+p\,\ b\,\ c]$, since a cyclic permutation of three forms preserves
the bracket. Likewise, $[s\,r\,p]=-[p\,r\,s]$ cancels the corresponding
term from $[p+q\,\ r\,\ s]$. The term $-[s\,c\,p]$ is
$[c\,s\,p]$, so the remaining terms give
\[
\begin{aligned}
&[a\,q\,p]+[b+s\,\ r-c\,\ p]
 +[a+p\,\ b\,\ c]+[p+q\,\ r\,\ s]\\
&\qquad=[a\,q\,p]+[b\,r\,p]+[c\,s\,p]
       +[a\,b\,c]+[q\,r\,s]\,.
\end{aligned}
\]
For each of the first three brackets on the right, the identity above
and $p=-\iota h$ give
\[
\begin{aligned}
[x_j+\iota y_j\,\ x_j-\iota y_j\,\ p]
&=-2\iota[x_j\,y_j\,-\iota h]\\
&=-2[x_j\,y_j\,h]
 =2[x_j\,h\,y_j]\,.
\end{aligned}
\]
In $[a\,b\,c]+[q\,r\,s]$, terms containing an odd number of $y$
coordinates have opposite coefficients and cancel. Terms containing no
$y$ coordinates have coefficient $2$, and terms containing two have
coefficient $2\iota^2=-2$. Reordering the forms in the latter terms gives
\[
\begin{aligned}
[a\,b\,c]+[q\,r\,s]
=2\bigl(&[x_1\,x_2\,x_3]-[x_1\,y_2\,y_3]\\
       &+[x_2\,y_1\,y_3]-[x_3\,y_1\,y_2]\bigr)\,.
\end{aligned}
\]
Comparison with \eqref{eq:octonion-three-form} now gives
\[
2\Phi=[a\,q\,p]+[b+s\,\ r-c\,\ p]
       +[a+p\,\ b\,\ c]+[p+q\,\ r\,\ s]\,.
\]
After division by $2$, the first bracket is
$-\iota[x_1\,y_1\,p]$. For any linear map $L$ with rows $f,g,\ell$,
ordinary matrix transposition gives
\[
\begin{aligned}
\langle L^{\mathsf T}(LX\times LY),Z\rangle_3
&=(LX\times LY)\cdot LZ\\
&=[f\,g\,\ell](X,Y,Z)\,.
\end{aligned}
\]
The four brackets are therefore the pairings with $Z$ of the four terms
on the right of \eqref{eq:octonion-leaf}. Their sum equals
$\langle\AltC_3(X,Y),Z\rangle_3$ for every $Z$. Pairing the difference
with each standard basis vector shows that all its coordinates vanish,
which proves \eqref{eq:octonion-leaf}.

For the arithmetic count, evaluate the input linear maps, the four
complex cross products, and the output reconstruction in order.

\emph{Input linear maps.} For each operand, forming $z$ costs $1$ complex
addition; forming $b,r,c,s$ costs $4$; forming $b+s,r-c$ costs $2$; and
forming $z\pm\iota y_1$ costs $2$ more. Thus the input maps require $9$
complex additions, or $18$ real additions, per operand. Multiplication
by $\iota$ only exchanges coordinates and changes a sign, so both input
maps together cost $36$ real additions.

\emph{Cross products.} Write
\[
(\Xi_j,\Psi_j,\Gamma_j):=L_jX\times L_jY\qquad(1\le j\le4)\,.
\]
Each complex cross product can be evaluated by three real cross products.
For $u=u_0+\iota u_1$ and $v=v_0+\iota v_1$, form
\[
\Delta=u_0\times v_0\,,\qquad \Omega=u_1\times v_1\,,\qquad
W=(u_0+u_1)\times(v_0+v_1)\,.
\]
Then $u\times v=(\Delta-\Omega)+\iota(W-\Delta-\Omega)$. The four
complex cross products therefore require twelve real cross products,
each using $6$ multiplications and $3$ subtractions, for a total of $72$
multiplications and $36$ additions. The sums $u_0+u_1$ and $v_0+v_1$
would cost $6$ additions per complex cross product, or $24$ in total.
The third coordinate of both $L_1$ and $L_2$ is $p$, so the sum of the
real and imaginary parts of $p$ is computed once and reused in each
operand. This saves $2$ additions, reducing the input-sum count to $22$.

For the $j$th complex cross product, denote its three real cross products
by $\Delta_j,\Omega_j,W_j$. The coordinates $\Gamma_1,\Gamma_2$ will be
combined directly when computing $\widehat h$, so only the other ten
complex output coordinates are reconstructed separately. Each uses $1$
subtraction for its real part and $2$ for its imaginary part, costing
$30$ additions in total. This stage therefore uses $72$ multiplications
and $36+22+30=88$ additions.

\emph{Output reconstruction.} Form the six complex intermediate
quantities
\[
\begin{gathered}
\Theta=\Xi_3+\Xi_4\,,\qquad U=\Xi_3-\Xi_4\,,\\
d_1=\Xi_2+\Psi_3\,,\qquad d_2=\Psi_2+\Psi_4\,,\\
e_1=\Gamma_3-\Psi_2\,,\qquad e_2=\Xi_2+\Gamma_4\,.
\end{gathered}
\]
Write the output in the same order as $X$ as
$(\widehat x_1,\widehat x_2,\widehat x_3,\widehat h,
\widehat y_1,\widehat y_2,\widehat y_3)$. The sparse transpose
reconstruction gives
\begin{equation}\label{eq:octonion-reconstruction}
\begin{aligned}
\widehat x_1&=\frac{\Theta}{2}-\iota \Xi_1\,,
&\widehat y_1&=\iota\left(\frac{U}{2}-\Psi_1\right)\,,\\
\widehat x_2&=\frac{d_1+d_2}{2}\,,
&\widehat y_2&=\frac{\iota(d_1-d_2)}{2}\,,\\
\widehat x_3&=\frac{e_1+e_2}{2}\,,
&\widehat y_3&=\frac{\iota(e_1-e_2)}{2}\,,\\
\widehat h&=-\frac{\iota(\Gamma_2+\Theta)}{2}-\Gamma_1\,.
\end{aligned}
\end{equation}
To obtain the last coordinate without reconstructing $\Gamma_1,\Gamma_2$,
let $p_j,\eta_j,\omega_j$ be the third coordinates of
$\Delta_j,\Omega_j,W_j$ for $j=1,2$. Then
\[
\Gamma_j=(p_j-\eta_j)+\iota(\omega_j-p_j-\eta_j)\,.
\]
Retain $\theta_R:=\ReC(\Theta/2)$ and
$\theta_I:=\ImC(\Theta/2)$ from the computation of $\widehat x_1$.
Taking real and imaginary parts of the formula for $\widehat h$ gives
\[
\begin{aligned}
\ReC\widehat h
&=\frac{\omega_2-p_2-\eta_2}{2}+\theta_I-p_1+\eta_1\,,\\
\ImC\widehat h
&=-\frac{p_2-\eta_2}{2}-\theta_R-\omega_1+p_1+\eta_1\,.
\end{aligned}
\]
Define
\[
e:=p_1+\frac{\eta_2}{2}\,,\qquad f:=\eta_1-\frac{p_2}{2}\,.
\]
Collecting the terms involving $p_1,\eta_1,p_2,\eta_2$ then yields
\[
\ReC\widehat h=f-e+\frac{\omega_2}{2}+\theta_I\,,\qquad
\ImC\widehat h=e+f-\omega_1-\theta_R\,.
\]
These formulas use the retained product coordinates directly and agree
with \eqref{eq:octonion-reconstruction}.

The six intermediate quantities require $6$ complex additions, and the
six output coordinates other than $\widehat h$ require another $6$,
giving $24$ real additions. Computing $e,f$ costs $2$ real additions,
and each coefficient of $\widehat h$ costs $3$, giving $8$ more. Halving
$\Theta$ and $U$ costs $4$ real halvings, and the four outputs
$\widehat x_2,\widehat y_2,\widehat x_3,\widehat y_3$ cost another $8$.
The last coordinate uses $3$ additional halvings for
$\eta_2/2,p_2/2,\omega_2/2$; the retained coefficients of $\Theta/2$
require none. Thus the three stages have costs
\[
\begin{array}{lrrr}
 & M&A&E\\
\text{Input linear maps}&0&36&0\\
\text{Four complex cross products}&72&88&0\\
\text{Output reconstruction}&0&32&15\\ \hline
\text{Total}&72&156&15
\end{array}
\]
and the total number of real operations is $243$.
\end{proof}

\section{Arithmetic counts}
\label{subsec:operation-counts}

To obtain the arithmetic counts for the recursion, first consider the
evaluation of $B_k^{\C}$ in real arithmetic by
\eqref{eq:complex-fusion}. The linear input programs form the factors
of the elementary products, and the linear output program reconstructs
the result from their values. Here a linear program consists of binary
additions or subtractions and fixed-scalar multiplications, with copying
and reuse allowed.

\begin{proposition}[Addition count for the nonrecursive
map]\label{prop:nonrecursive-additions}
Let $k\ge2$ and $m=2^{k-1}-1$. Consider the evaluation of $B_k^{\C}$ by
\eqref{eq:complex-fusion}, using the six coordinate products
in \eqref{eq:nonrecursive-map} for each real evaluation of $B_k$.
For the given input coordinates and
prescribed product formulas, with the product values treated as
independent inputs to the output program, the two linear input programs
and the linear output program require at least $22m+2$ real additions
in total, even when fixed real scalar multiplications are free.
The specified evaluation attains this bound and uses $18m$
input-dependent real products.
\end{proposition}

\begin{proof}
Put $q=2m+1$, the number of coordinates in each input and output of
$B_k$. For each $1\le i\le m$, the prescribed elementary products in
one real evaluation are
\[
\delta s_i\,,\qquad\gamma t_i\,,\qquad
\delta u_i\,,\qquad\gamma v_i\,,\qquad
s_iv_i\,,\qquad t_iu_i\,.
\]
The parameters and pairing weights in the standard tower are signs,
so they do not require additional arithmetic operations. To apply
\eqref{eq:complex-fusion}, write the complex inputs as
\[
\begin{aligned}
(u,\gamma,s)&=(u_0,\gamma_0,s_0)
              +\iota(u_1,\gamma_1,s_1)\,,\\
(v,\delta,t)&=(v_0,\delta_0,t_0)
              +\iota(v_1,\delta_1,t_1)\,,
\end{aligned}
\]
where all quantities on the right have real coordinates. The evaluations
$\Delta$ and $\Omega$ use the real and imaginary parts directly. The
evaluation $W$ instead requires the following coordinate sums as its
input factors:
\[
\begin{array}{lll}
u_{0,i}+u_{1,i}\,,&s_{0,i}+s_{1,i}\,,&\gamma_0+\gamma_1\,,\\
v_{0,i}+v_{1,i}\,,&t_{0,i}+t_{1,i}\,,&\delta_0+\delta_1\,.
\end{array}
\]
Here $i$ ranges from $1$ to $m$ in the vector sums, while each scalar
sum is formed once. Every listed sum occurs in at least one prescribed
product because $m\ge1$. There are therefore $q$ required sums for
each operand, and computing them directly uses $2q$ additions.
Each real evaluation
uses $6m$ multiplications, $2m$ subtractions for its vector outputs, and
$2m-1$ additions and subtractions for its scalar output. Reconstructing
each complex output from the three real evaluations uses $3$ additions,
so the specified evaluation has multiplication count $18m$ and addition
count
\[
2q+3(4m-1)+3q=22m+2\,.
\]

For the lower bound, the prescribed products require each input program
to produce all of its $q$ listed sums, up to nonzero fixed scaling. For
each operand, its real and imaginary input coordinates are independent
variables. The required sums involve different pairs of these variables,
so they are distinct up to nonzero scalar multiplication, and none is a
scalar multiple of an input coordinate. Each required sum must therefore
first appear at an addition or subtraction step. Such a step creates at
most one new linear form up to scaling, while multiplication by a fixed
scalar creates none. Thus each input program needs at least $q$ additions.
The two input programs together require at least $2q$ additions,
independently of the output program.

For a real matrix $J$, let $a(J)$ denote the minimum number of additions
and subtractions needed to evaluate its linear map when fixed scalar
multiplications are free. Let $W_B$ be the $2q\times18m$ matrix mapping
the independent elementary product values to the real and imaginary
output coordinates. The products contributing to each vector coordinate
of a real evaluation of $B_k$ are its two scalar--coordinate products.
The scalar coordinate instead receives the $2m$ products in the pairing
difference. No product contributes to another coordinate of that real
evaluation.

For the $j$th output coordinate, let $r_j$ be the number of these
products and let $\varepsilon_j$ be the row vector of their signs.
Thus $r_j=2$ for each of the $2m$ vector coordinates and $r_j=2m$ for
the scalar coordinate, with every entry of $\varepsilon_j$ equal to
$1$ or $-1$. In particular, $\sum_j r_j=6m$. Ordering the product
columns by the evaluations $\Delta$, $\Omega$, and $W$, the block of
$W_B$ for this real-imaginary output pair is
\[
W_{B,j}=
\begin{pmatrix}
\varepsilon_j&-\varepsilon_j&0\\
-\varepsilon_j&-\varepsilon_j&\varepsilon_j
\end{pmatrix}\,.
\]
The zero block has $r_j$ entries. This follows directly from the
reconstruction $(\Delta-\Omega)+\iota(W-\Delta-\Omega)$, applied
coordinate by coordinate. After reordering rows and columns, $W_B$
consists of these independent blocks. Every column is nonzero, and both
rows of every block are nonzero, so $W_B$ has no zero rows or columns.

Let $\sigma_j,\chi_j$ be the two inputs to the transposed block. Its
outputs are, up to their signs and repetitions,
\[
\sigma_j-\chi_j\,,\qquad \sigma_j+\chi_j\,,\qquad\chi_j\,.
\]
Across the $q$ pairs, the $2q$ sums and differences are distinct up to
nonzero scaling and are not scalar multiples of input coordinates.
The same new-form argument used for the input programs therefore gives
$a(W_B^{\mathsf T})\ge2q$. Computing these sums and differences directly
uses $2q$ additions, after which all other outputs are copies or sign
changes. Hence
\[
a(W_B^{\mathsf T})=2q\,.
\]

The transposition principle of
\href{https://doi.org/10.1016/0196-6774(88)90026-0}{Kaminski, Kirkpatrick,
and Bshouty (1988)} states that, for a matrix $J$ without zero rows or
columns, with $d_{\mathrm{out}}$ rows and $d_{\mathrm{in}}$ columns, a
linear program can be transposed with an addition-count change of
$d_{\mathrm{out}}-d_{\mathrm{in}}$. Applied to optimal programs in both
directions, this gives
\[
\begin{aligned}
a(J^{\mathsf T})&\le a(J)+d_{\mathrm{out}}-d_{\mathrm{in}}\,,\\
a(J)&\le a(J^{\mathsf T})+d_{\mathrm{in}}-d_{\mathrm{out}}\,,
\end{aligned}
\]
and therefore $a(J^{\mathsf T})=a(J)+d_{\mathrm{out}}-d_{\mathrm{in}}$.
Substituting $d_{\mathrm{out}}=2q$ and $d_{\mathrm{in}}=18m$ yields
\[
a(W_B)=a(W_B^{\mathsf T})+18m-2q=18m\,.
\]
Together with the $2q$ additions required by the input programs, this
gives the lower bound $18m+2q=22m+2$. The evaluation counted at the
beginning attains it, completing the proof.
\end{proof}

\begin{proposition}[Exact arithmetic counts]\label{prop:exact-counts}
The recursive schedule of Section~\ref{sec:notation}, with the octonion base
case of
Proposition~\ref{prop:octonion-leaf} and the one-call real-input root of
Corollary~\ref{cor:real-root}, multiplies two elements of the standard
real algebra $A_n$, $n\ge4$, $N=2^n$, using
\begin{equation}\label{eq:ql-arithmetic-counts}
\begin{aligned}
M_{\mathrm{QL}}&=\frac{(36n-69)N}{8}+10\,,
&A_{\mathrm{QL}}&=\frac{(76n-196)N}{8}+26\,,\\
E_{\mathrm{QL}}&=\frac{(16n-53)N}{16}+4\,,
&T_{\mathrm{QL}}&=\frac{(240n-583)N}{16}+40\,.
\end{aligned}
\end{equation}
For $n=1,2,3$, the respective counts $(M,A,E)$ are $(4,2,0)$,
$(16,12,0)$, and $(58,62,0)$.
\end{proposition}

\begin{proof}
Let $M_k,P_k,H_k$ denote the numbers of input-dependent real products,
real additions and subtractions, and real halvings, respectively, in one
evaluation of the complex kernel $\AltC_k$, including its recursive
calls. Proposition~\ref{prop:octonion-leaf} gives the base counts
\[
(M_3,P_3,H_3)=(72,156,15)\,.
\]

\emph{Complex nodes.} Fix $k\ge4$ and put $m=2^{k-1}-1$, so that
$u,v,s,t$ have $m$ complex coordinates each. By
Proposition~\ref{prop:nonrecursive-additions}, the nonrecursive map uses $18m$
real
multiplications and $22m+2$ additions. Each coordinate of
$u\pm\iota s$ requires $2$ real additions for each sign, and the same
holds for $v\pm\iota t$, giving $8m$ additions for the transformed
inputs. Each recursive output has $m$ complex coordinates, so forming
$Z_++Z_-$ and $Z_+-Z_-$ uses $4m$ real additions. Combining them with
the two vector outputs of the nonrecursive map costs another $4m$ additions.
The scalar output $\lambda$ is inserted in the otherwise zero unit
coordinate of the second component and requires no addition.

Halving the sum and difference of the recursive outputs costs $4m$ real
halvings. Multiplication by $\iota$ only exchanges coordinates and
changes signs, so it has no additional arithmetic cost. The work outside
the recursive calls is therefore
\[
\begin{array}{lrrr}
 & M&A&E\\
\text{Nonrecursive map}&18m&22m+2&0\\
\text{Transformed inputs}&0&8m&0\\
\text{Recursive sums and differences}&0&4m&0\\
\text{Vector recombination}&0&4m&0\\
\text{Halving}&0&0&4m\\ \hline
\text{Total}&18m&38m+2&4m
\end{array}
\]
Including both recursive calls gives
\[
\begin{aligned}
M_k&=2M_{k-1}+18(2^{k-1}-1)\,,\\
P_k&=2P_{k-1}+38(2^{k-1}-1)+2\,,\\
H_k&=2H_{k-1}+4(2^{k-1}-1)\,.
\end{aligned}
\]

At level $j$ of a kernel evaluation beginning at level $k$, there are
$2^{k-j}$ nodes, and the recursion ends with $2^{k-3}$ octonion leaves.
Unrolling the recurrences therefore gives
\[
\begin{aligned}
M_k&=72\cdot2^{k-3}
 +18\sum_{j=4}^k2^{k-j}(2^{j-1}-1)\,,\\
P_k&=156\cdot2^{k-3}
 +38\sum_{j=4}^k2^{k-j}(2^{j-1}-1)
 +2\sum_{j=4}^k2^{k-j}\,,\\
H_k&=15\cdot2^{k-3}
 +4\sum_{j=4}^k2^{k-j}(2^{j-1}-1)\,.
\end{aligned}
\]
For $k\ge3$, with empty sums interpreted as zero,
\[
\begin{aligned}
\sum_{j=4}^k2^{k-j}&=2^{k-3}-1\,,\\
\sum_{j=4}^k2^{k-j}(2^{j-1}-1)
&=(k-3)2^{k-1}-2^{k-3}+1\,.
\end{aligned}
\]
Substitution and collection of terms now yield
\[
\begin{aligned}
M_k&=(36k-81)2^{k-2}+18\,,\\
P_k&=(76k-168)2^{k-2}+36\,,\\
H_k&=(16k-37)2^{k-3}+4\,.
\end{aligned}
\]

\emph{Real root.} For $n\ge4$, $\AltR_n$ forms $u+\iota s$ and
$v+\iota t$ by interleaving real coordinates, at no arithmetic cost, and
makes one call to $\AltC_{n-1}$. With $m=2^{n-1}-1$, the real nonrecursive map
uses $6m$ multiplications and $4m-1$ additions. Combining its two
vector outputs with the real and imaginary parts of the recursive
output costs another $2m$ additions. The work at the root therefore
adds $6m$ multiplications and $6m-1=3N-7$ additions, with no halvings.

\emph{Reconstruction.} In \eqref{eq:scalar--imaginary}, the scalar
product $\alpha\beta$ requires $1$ multiplication, the pairing requires
$N-1$, and each scalar--vector product requires $N-1$. Their total is
\[
1+(N-1)+2(N-1)=3N-2
\]
input-dependent real multiplications. Summing the terms of the pairing
costs $N-2$ additions, and subtracting it from $\alpha\beta$ costs $1$
more. Forming $\alpha Y+\beta X$ costs $N-1$ additions, and adding
$X\times_nY$ costs another $N-1$, giving
\[
(N-2)+1+(N-1)+(N-1)=3N-3
\]
additions. The standard pairing weights are all $1$, so reconstruction
requires no additional fixed-scalar multiplications.

All charged fixed-scalar operations are consequently halvings in the
complex kernel. Summing the kernel call, real root, and reconstruction
costs gives
\[
\begin{aligned}
M_{\mathrm{QL}}
&=M_{n-1}+6(2^{n-1}-1)+3N-2
 =M_{n-1}+6N-8\,,\\
A_{\mathrm{QL}}
&=P_{n-1}+(3N-7)+(3N-3)
 =P_{n-1}+6N-10\,,\\
E_{\mathrm{QL}}&=H_{n-1}\,.
\end{aligned}
\]
Substituting $k=n-1$ into the kernel formulas gives the first three
counts in \eqref{eq:ql-arithmetic-counts}. Adding them gives
\[
T_{\mathrm{QL}}=M_{\mathrm{QL}}+A_{\mathrm{QL}}+E_{\mathrm{QL}}
=\frac{(240n-583)N}{16}+40\,.
\]

\emph{Small dimensions.} At $n=1$, the alternating product vanishes,
and the full product is
\[
(\alpha+x_1e_1)(\beta+y_1e_1)
=(\alpha\beta-x_1y_1)+(\alpha y_1+\beta x_1)e_1\,.
\]
Evaluating the four displayed scalar products uses $4$ multiplications,
and the two output coordinates use $2$ additions, giving $(4,2,0)$.
At $n=2$, the real alternating product is the ordinary three-dimensional
cross product, with $6$ multiplications and $3$ subtractions. Full
reconstruction adds $10$ multiplications and $9$ additions, so the
counts are $(16,12,0)$.

At $n=3$, the complex call is $\AltC_2$, whose nonrecursive map has $m=1$ and
$q=3$. Its three real evaluations use $18$ multiplications and $9$
additions. The input sums use $6$ additions, and the complex output
reconstruction uses $9$, giving kernel counts $(18,24,0)$. The real
root has $m=3$ and adds $18$ multiplications and $17$ additions, while
full reconstruction adds $22$ multiplications and $21$ additions. Thus
the full counts are
\[
(18+18+22,\;24+17+21,\;0)=(58,62,0)\,.
\]
\end{proof}

\begin{corollary}[Real bilinear bound]\label{cor:real-rank}
For any fixed nonzero real doubling parameters, use the real-input
schedule of Section~\ref{subsec:real-specialization}, with the complex
octonion base case of Proposition~\ref{prop:octonion-leaf} transported by
fixed linear coordinate changes. Let $C_n$ denote its number of
input-dependent real products. Then
\begin{equation}\label{eq:Cn}
C_n\le(36n-93)2^{n-3}+12\qquad(n\ge4)\,.
\end{equation}
Consequently,
\begin{equation}\label{eq:rank-bound}
\Rk_{\R}(T_{A_n})\le(36n-69)2^{n-3}+10\qquad(n\ge4)\,.
\end{equation}
For $n=1,2,3$, this schedule gives the nonoptimal upper bounds $4$,
$16$, and $58$, respectively.
\end{corollary}

\begin{proof}
First transport the complex octonion formula to the given parameters.
For $1\le k\le3$, choose $\tau_k\in\C^\times$ with
$\tau_k^2=-\mu_k$, define $D_0$ to be the identity on $\C$, and put
\[
D_k(a,b):=(D_{k-1}a,\tau_kD_{k-1}b)\,.
\]
Inductively, $D_k$ is an isomorphism from the complexification of the
parameterized level-$k$ algebra to the complexification of the standard
one. To check this, suppose that $D_{k-1}$ is such an isomorphism and
commutes with the involution, and write $\widetilde z=D_{k-1}z$ for
$z=a,b,c,d$. The involution is $\C$-linear, so it fixes $\tau_k$.
Multiplication of the images in the standard algebra therefore gives
\[
\begin{aligned}
&(\widetilde a,\tau_k\widetilde b)
 (\widetilde c,\tau_k\widetilde d)\\
&\quad=
 \bigl(\widetilde a\widetilde c
       -\tau_k^2\overline{\widetilde d}\,\widetilde b,\;
       \tau_k(\widetilde d\widetilde a
                   +\widetilde b\overline{\widetilde c})\bigr)\\
&\quad=D_k\bigl(ac+\mu_k\bar d b,\;da+b\bar c\bigr)\,,
\end{aligned}
\]
where the last line uses $\tau_k^2=-\mu_k$ and the induction
hypothesis. This is the image of the parameterized product in
\eqref{eq:CD-product}. Moreover,
\[
\overline{D_k(a,b)}
=(D_{k-1}\bar a,-\tau_kD_{k-1}b)
=D_k(\bar a,-b)\,.
\]
Since $\tau_k\ne0$ and $D_{k-1}$ is invertible, $D_k$ is invertible
as well. Starting from $D_0$, this proves the asserted isomorphisms and
their compatibility with the involution. The maps are diagonal in the
recursive standard coordinates and preserve the imaginary subspaces.

Write $\AltC_3^{(\mu)}$ and $\AltC_3^{\mathrm{std}}$ for the
parameterized and standard alternating products. Preservation of
multiplication and of the involution gives
\[
\AltC_3^{(\mu)}(X,Y)
=D_3^{-1}\AltC_3^{\mathrm{std}}(D_3X,D_3Y)\,.
\]
Upon writing each complex coordinate as two real coordinates, $D_3$ and
$D_3^{-1}$ are fixed real-linear maps. They change the input linear
forms and the output recombination but introduce no input-dependent
products. Proposition~\ref{prop:octonion-leaf} therefore supplies a
complex alternating-product evaluation with $72$ input-dependent real
products for any nonzero real parameters. In that evaluation, the input
maps $L_j$ and the sums used to evaluate the complex cross products act
separately on the two operands. Each input-dependent real multiplication
consequently multiplies a real-linear form in the first operand by a
real-linear form in the second, and the output reconstruction is
real-linear in these products. Composing the input maps with $D_3$ and
the output map with
$D_3^{-1}$ preserves this property. The fixed scalings in the
coordinate changes can affect the total arithmetic count, so the exact
total in Proposition~\ref{prop:exact-counts} remains a statement about
the standard tower.

For the complex kernel, the level-two formula uses $M_2=18$ products and
the transported level-three formula uses $M_3=72$. At higher levels, the
parameters and pairing weights contribute only fixed linear operations,
while the nonrecursive map uses $18(2^{k-1}-1)$ input-dependent products.
Consequently, the same product recurrence as in
Proposition~\ref{prop:exact-counts} applies, giving
\[
M_k=(36k-81)2^{k-2}+18\qquad(k\ge3)\,.
\]

The same bilinear property holds throughout the complex recursion. More
precisely, in real coordinates, each input-dependent multiplication is
the product of a fixed real-linear form in the first kernel input and a
fixed real-linear form in the second, and every output coordinate is a
fixed real-linear combination of the resulting products. At level $2$,
this follows from the elementary products of $B_2$ and the fusion formula
\eqref{eq:complex-fusion}. The transported octonion evaluation establishes
the assertion at level $3$. For $k\ge4$, assume it at level $k-1$.
Multiplication by the fixed coefficient $\rho_k$ and formation of the
sums and differences in the recursive inputs are real-linear operations
on each operand separately. Substituting these transformed inputs into
the linear forms from the induction hypothesis therefore gives
real-linear forms in the original kernel inputs. The nonrecursive map has the
same property: each real evaluation of $B_k$ uses elementary products
with one factor from each operand, and the fusion formula forms its
inputs separately and recombines its outputs linearly. Finally,
\eqref{eq:two-call-formula} uses only fixed real-linear combinations of
the outputs of the nonrecursive map and the recursive calls. This proves the
assertion by induction.

For real inputs, the nonrecursive map at level $n$ uses
$6(2^{n-1}-1)$ products. If $\mu_n<0$, one complex call at the preceding
level suffices by Corollary~\ref{cor:real-root}. If $\mu_n>0$, the
transformed inputs are real and two real-input calls suffice instead.
The real base formulas give $C_1=0$ and $C_2=6$. At level $3$, the
nonrecursive map uses $6\cdot3=18$ products, so the negative-parameter
case gives $C_3\le M_2+18=36$, while the positive-parameter case gives
$C_3\le2C_2+18=30$. Hence $C_3\le36$ for either sign.

The real-input schedule also has the bilinear property just proved. Its
level-one and level-two formulas have this property directly. At a
higher level with $\mu_n<0$, it follows from the property of the complex
kernel, since the transformed inputs are real-linear in the separate
operands and both $\ReC$ and $\ImC$ are real-linear maps. When
$\mu_n>0$, the transformed inputs are real, so induction applies to the
two real-input calls. In both cases, the real nonrecursive map pairs factors
from the separate operands, and the remaining output operations are
real-linear.

Put $U_n=(36n-93)2^{n-3}+12$ for $n\ge4$. At level $4$, the nonrecursive map
uses $6\cdot7=42$ products. The negative-parameter case gives
$C_4\le M_3+42=114$, and the positive-parameter case gives
$C_4\le2C_3+42\le114$. Since $U_4=114$, the bound holds at level $4$.

Now let $n\ge5$ and suppose that $C_{n-1}\le U_{n-1}$. If $\mu_n<0$,
substitution of the complex-kernel formula yields
\[
\begin{aligned}
C_n&\le M_{n-1}+6(2^{n-1}-1)\\
&=(36n-117)2^{n-3}+18+24\cdot2^{n-3}-6\\
&=(36n-93)2^{n-3}+12=U_n\,.
\end{aligned}
\]
If $\mu_n>0$, the two real-input calls and the induction hypothesis give
\[
\begin{aligned}
C_n&\le2U_{n-1}+6(2^{n-1}-1)\\
&=(36n-105)2^{n-3}+18\\
&=U_n-\bigl(12\cdot2^{n-3}-6\bigr)\le U_n\,,
\end{aligned}
\]
since $12\cdot2^{n-3}-6\ge0$ for $n\ge5$. This proves
\eqref{eq:Cn} for every choice of signs of the parameters.

Full reconstruction in \eqref{eq:scalar--imaginary} adds $3N-2$
input-dependent products. Each of these also multiplies a real-linear
form in $x$ by a real-linear form in $y$: they are the scalar product
$\alpha\beta$, the weighted coordinate products in the pairing, and the
coordinate products in $\alpha Y$ and $\beta X$. Since extracting the
scalar and imaginary coordinates is real-linear, the products in the
alternating-product schedule likewise have this form in the full inputs
$x$ and $y$. Collecting the fixed output coefficients of all counted
products therefore gives a decomposition
\[
xy=\sum_{j=1}^r\alpha_j(x)\beta_j(y)v_j\,,\qquad
r\le C_n+3N-2\,,
\]
where $\alpha_j,\beta_j$ are real-linear forms and each $v_j\in A_n$
is fixed. By the definition of real tensor rank, this proves
$\Rk_{\R}(T_{A_n})\le C_n+3N-2$. For $n\ge4$,
\[
\begin{aligned}
C_n+3N-2
&\le(36n-93)2^{n-3}+12+24\cdot2^{n-3}-2\\
&=(36n-69)2^{n-3}+10\,,
\end{aligned}
\]
which gives \eqref{eq:rank-bound}. At $n=1,2,3$, the corresponding
counts are $C_1+4=4$, $C_2+10=16$, and $C_3+22\le58$, respectively.
\end{proof}

For direct multiplication, each Cayley--Dickson step evaluates four
products of dimension $N/2$ and combines their coordinates using $N$
additions and subtractions. If $M_{\mathrm{direct}}(n)$ and
$A_{\mathrm{direct}}(n)$ denote its operation counts at level $n$, then
\[
\begin{aligned}
M_{\mathrm{direct}}(n)&=4M_{\mathrm{direct}}(n-1)\,,\\
A_{\mathrm{direct}}(n)&=4A_{\mathrm{direct}}(n-1)+2^n\,,
\end{aligned}
\]
with $M_{\mathrm{direct}}(0)=1$ and $A_{\mathrm{direct}}(0)=0$.
Consequently,
\[
M_{\mathrm{direct}}=4^n=N^2\,,\qquad
A_{\mathrm{direct}}
=\sum_{j=1}^n4^{n-j}2^j=N^2-N\,.
\]
The only fixed coefficients in the standard product are signs, so
$E_{\mathrm{direct}}=0$ and
\[
T_{\mathrm{direct}}=2N^2-N\,.
\]

For $N\ge2$, the uniform Cariow--Cariowa (CC) method uses a Hadamard
factorization
with a sparse correction of the basis-product signs. The counts in
\href{https://doi.org/10.15199/48.2015.02.09}{Cariow and Cariowa
(2015, p.~38)} can be recovered as follows, with the power-of-two
scalings included under the present convention.

In the standard basis, induction in \eqref{eq:CD-product} gives
\[
e_i e_j=\sigma_n(i,j)e_{i\oplus j}\,,\qquad
\sigma_n(i,j)\in\{1,-1\}\,,
\]
where $\oplus$ is bitwise exclusive-or on the binary indices. The basis
unit gives $\sigma_n(0,j)=\sigma_n(j,0)=1$, and the imaginary basis
elements satisfy
\[
\sigma_n(i,i)=-1\quad(i>0)\,,\qquad
\sigma_n(i,j)=-\sigma_n(j,i)\quad(i,j>0,\ i\ne j)\,.
\]
Indeed, these last identities follow from
$e_i e_j+e_j e_i=-2\langle e_i,e_j\rangle_n$ and the standard pairing.
Put $\epsilon_0=-1$ and $\epsilon_k=1$ for $k>0$, and define
\[
\mathcal I_k:=
\{i:\epsilon_k\sigma_n(i,i\oplus k)=-1\}\,.
\]
The unsigned convolution coordinates
$c_k:=\sum_i x_i y_{i\oplus k}$ then recover the product coordinates by
\[
(xy)_k=\epsilon_k\left(c_k-
 2\sum_{i\in\mathcal I_k}x_i y_{i\oplus k}\right)\,.
\]
Thus the correction replaces exactly the terms whose signs differ from
the unsigned convolution, followed by the free sign change $\epsilon_k$.

For $k=0$, the correction set is $\mathcal I_0=\{0\}$. For $k>0$, the
indices form $N/2$ disjoint pairs $\{i,i\oplus k\}$. The pair
$\{0,k\}$ has positive basis-product signs and contributes no
correction. In every other pair, both indices are nonzero and their
basis products have opposite signs. Exactly one index in that pair
therefore belongs to $\mathcal I_k$, giving
\[
|\mathcal I_0|=1\,,\qquad
|\mathcal I_k|=N/2-1\quad(k>0)\,.
\]
The total number of correction products is consequently
\[
\sum_{k=0}^{N-1}|\mathcal I_k|
=1+(N-1)(N/2-1)=\frac{N(N-3)+4}{2}\,.
\]

To evaluate the convolution, let $\mathcal H_N$ be the Walsh--Hadamard matrix
with entries $(\mathcal H_N)_{a,b}=(-1)^{a\cdot b}$, where $a\cdot b$ is the
binary digit dot product modulo $2$. It is symmetric and satisfies
$\mathcal H_N^2=NI_N$, where $I_N$ is the $N\times N$ identity matrix.
Consequently,
\[
c=\frac1N\mathcal H_N
 \bigl((\mathcal H_Nx)\odot(\mathcal H_Ny)\bigr)\,,
\]
where $\odot$ denotes coordinatewise multiplication. In fact, expansion
of the right-hand side gives coefficient
$N^{-1}\sum_a(-1)^{a\cdot(i\oplus j\oplus k)}$ on $x_i y_j$ in
coordinate $k$, which is $1$ when $i\oplus j=k$ and $0$ otherwise.
Each transform has $n$ stages of $N/2$ butterflies, each using $2$
additions or subtractions, and therefore costs $nN$ additions. The three
transforms use $3nN$ additions, the coordinatewise products use $N$
multiplications, and normalization uses $N$ fixed-scalar multiplications.

For a nonempty correction set of size $r$, summing its products costs
$r-1$ additions and subtracting the doubled sum costs $1$ more. An
empty set requires neither operation. Hence the corrections require
$\sum_k|\mathcal I_k|$ additions as well as that many multiplications.
It follows that
\[
M_{\mathrm{CC}}=\frac{N(N-1)}2+2\,,\qquad
A_{\mathrm{CC}}=3nN+\frac{N(N-3)+4}{2}\,.
\]
The normalizations and correction doublings are omitted in the published
count because they are power-of-two scalings. When $N\ge4$, every
correction set is nonempty, so the schedule charges $N$ normalizations
and $N$ doublings, giving $E_{\mathrm{CC}}=2N$. At $N=2$, only
$\mathcal I_0=\{0\}$ is nonempty. The outputs are then
\[
(xy)_0=-\bigl(c_0-2x_0y_0\bigr)\,,\qquad (xy)_1=c_1\,,
\]
so normalization and the only nonzero correction coordinate give
$E_{\mathrm{CC}}=3$. In particular, the full counts at $N=2$ are
$(3,7,3)$.

For completeness, the comparison thresholds can be checked directly
from these formulas. For $n\ge4$ and $N=2^n$, the normalized differences
are
\[
\begin{aligned}
\frac{M_{\mathrm{direct}}-M_{\mathrm{QL}}}{N}
 &=N-\frac{36n-69}{8}-\frac{10}{N}\,,\\
\frac{M_{\mathrm{CC}}-M_{\mathrm{QL}}}{N}
 &=\frac N2-\frac{36n-65}{8}-\frac8N\,,\\
\frac{T_{\mathrm{direct}}-T_{\mathrm{QL}}}{N}
 &=2N-\frac{240n-567}{16}-\frac{40}{N}\,,\\
\frac{T_{\mathrm{CC}}-T_{\mathrm{QL}}}{N}
 &=N-\frac{192n-583}{16}-\frac{36}{N}\,,\\
\frac{M_{\mathrm{CC}}+A_{\mathrm{CC}}-T_{\mathrm{QL}}}{N}
 &=N-\frac{192n-551}{16}-\frac{36}{N}\,.
\end{aligned}
\]
All these expressions increase with $n$ for $n\ge4$. Their increments
from $n$ to $n+1$ are bounded below, respectively, by
$N-9/2$, $N/2-9/2$, $2N-15$, and $N-12$ for each of the last two;
these are positive for $N\ge16$. At $n=4$, the five differences are
\[
6\,,\quad -\frac{19}{8}\,,\quad \frac{79}{16}\,,\quad
\frac{35}{16}\,,\quad \frac3{16}\,.
\]
The second becomes $11/8>0$ at $n=5$. Together with the small-dimensional
counts, this shows that the quasilinear schedule uses fewer
input-dependent multiplications than direct multiplication from $N=8$
onward and than CC from $N=32$ onward. Its total equals direct
multiplication at $N=2,4,8$ and is smaller from $N=16$ onward, while it
is smaller than CC for every $N\ge2$. At $N=16$, the quasilinear total
is $417$, compared with $496$ for direct multiplication and $452$ for
CC. The last normalized difference also proves the advantage over CC
for $N\ge16$ when fixed-scalar costs are omitted, including
$M_{\mathrm{CC}}+A_{\mathrm{CC}}=420>417$ at $N=16$.


\begin{thebibliography}{Lem26}
\bibitem[Lem26]{Lemley2026}
Harrison Lemley.
\emph{Quasilinear multiplication in the Cayley--Dickson tower}.
\href{https://arxiv.org/abs/2609.11588}{arXiv:2609.11588}, 2026.
\end{thebibliography}

\end{document}

